\begin{frame}{자주 묻는 질문 (3.1)}
  \begin{columns}[T]
    \begin{column}{0.5\textwidth}
      \textbf{Q1. 생성 vs 판별, 뭐가 더 좋은가?}
      \begin{itemize}
        \item \textbf{데이터 적을 때}: 생성적(GDA, 나이브베이즈) ---
          ``클래스가 데이터를 어떻게 만드나''라는 \textbf{강한 가정}을
          깔고 들어가 적은 데이터로도 안정적.
        \item \textbf{데이터 많을 때}: 판별적(로지스틱회귀) ---
          가정 없이 직접 경계 학습 $\Rightarrow$ 가정 틀림 위험 없음.
        \item 3.2절에서 GDA와 로지스틱회귀가 \textbf{같은 형태의
          결정 경계}로 수렴함을 직접 확인.
        \item 더 깊이: Stanford CS229 노트.
      \end{itemize}
    \end{column}
    \begin{column}{0.5\textwidth}
      \textbf{Q2. $P(x)$ 를 왜 무시해도 되나?}
      \begin{itemize}
        \item 분류에서는 \textbf{고정된 $x$ 하나}를 놓고
          ``$y=0$ 이 맞을까, $y=1$ 이 맞을까''만 비교.
        \item $P(x)$ 는 어떤 클래스를 비교하든 \textbf{항상 같은 값}
          (분모가 같으면 분자만 비교해도 대소관계 유지).
        \item 그래서 \textbf{분자 $P(x|y)P(y)$ 만 최대화}하는 클래스를
          고르면 충분.
      \end{itemize}
    \end{column}
  \end{columns}
\end{frame}
